% ============================================================
% Chapter II — Literature Review
% ============================================================
\section{Literature Review}

\subsection{Theoretical Foundation}

The theoretical basis for this study draws on two frameworks: the physics of
tire degradation and fuel mass effects in motorsport, and the statistical
theory of multinomial classification with interaction modeling.

From a physical standpoint, tire performance in Formula One follows a
degradation curve governed by thermal cycling, mechanical wear, and compound
properties \parencite{baldi2022}. Simultaneously, the car's mass decreases
linearly as fuel burns off, directly affecting lap time through reduced
inertia and improved cornering grip \parencite{oldknow2022fuel}. These two
forces---tire degradation (increasing lap time) and fuel burn-off (decreasing
lap time)---operate in opposition throughout a race stint. Their interaction
represents the net physical state of the car at any given lap.

From a statistical standpoint, multinomial logistic regression provides the
appropriate inferential framework when the response variable has three or more
unordered categories \parencite{agresti2013}. When predictors are physically
related (e.g., tire age and fuel mass both correlate with race distance),
mean-centering before creating interaction products prevents structural
multicollinearity and preserves the interpretability of individual
coefficients \parencite{hosmer2013}. The Box's M test serves as a gate for
method selection: if covariance homogeneity holds, discriminant analysis is
appropriate; if it fails, multinomial logistic regression is the correct
choice \parencite{tabachnick2019}.

\subsection{Review of Existing Studies}

\subsubsection{Racing Analytics and Prediction}

\textcite{fister2021} conducted a systematic review of computational
intelligence methods in motorsport, identifying lap time prediction, tire
strategy optimization, and race outcome modeling as the three dominant
research streams. Their survey noted that most studies rely on aggregated
season-level features rather than within-race telemetry, leaving a gap in
modeling the physical dynamics that determine finishing order.

\textcite{kuhn2022f1} applied gradient-boosted trees and neural networks to
predict finishing positions across the 2016--2020 F1 seasons using driver
characteristics, constructor budgets, and circuit features. While their model
achieved 71\% top-5 accuracy, it treated each race as independent and did not
incorporate tire degradation or fuel load dynamics.

\textcite{song2020} used deep learning on historical F1 data (1950--2019) but
acknowledged that the model's inability to capture within-race state changes
limited its predictive power for individual Grand Prix outcomes.

\subsubsection{Tire Degradation Modeling}

\textcite{perdomo2021tire} developed a physics-informed neural network to
model Pirelli tire wear using thermal and mechanical telemetry from F1
sessions. Their model demonstrated that tire surface temperature, lateral
forces, and cumulative laps on a compound jointly determine the degradation
curve---supporting the use of tire age as a proxy variable in outcome models.

\textcite{baldi2022} modeled tire performance degradation in motorsport using
nonlinear regression, showing that degradation follows an exponential rather
than linear trajectory after a critical wear threshold. This informed our
decision to include tire age as a continuous predictor while acknowledging
the linear approximation may understate late-stint degradation effects.

\subsubsection{Fuel Load Effects}

\textcite{oldknow2022fuel} analyzed the competitive impact of the 2019 F1 fuel
load regulations, demonstrating that the 110\,kg fuel limit introduced
strategic variation---teams operating below the limit on Sundays gained lap
time at the cost of fuel management risk. Their analysis supports the use of a
linear fuel burn-off proxy as a reasonable first-order approximation of the
mass--performance relationship.

\subsubsection{Multinomial Logistic Regression in Sports}

\textcite{agresti2013} provides the theoretical foundation for multinomial
logistic regression as a generalization of binary logistic regression to
nominal response variables with more than two categories.

\textcite{hosmer2013} established diagnostic procedures for logistic
regression models, including the use of VIF for multicollinearity detection
and the importance of interaction terms when theoretically justified.

\subsubsection{Cross-Validation Strategies}

\textcite{varoquaux2019} evaluated cross-validation strategies for grouped
data and demonstrated that standard $k$-fold CV produces optimistically biased
performance estimates when observations within groups are correlated.
\textcite{wong2023} applied leave-one-group-out cross-validation to sports
analytics, confirming LOGO-CV provides more conservative and realistic
estimates of generalizability.

\subsubsection{Class Imbalance}

\textcite{he2009learning} surveyed learning from imbalanced datasets and
identified balanced class weighting as a baseline correction strategy that
adjusts the loss function to penalize misclassifications of the minority class
more heavily.

\subsubsection{Philippine Studies}

\textcite{dalisay2022} applied logistic regression to predict game outcomes
in the Philippine Basketball Association using box-score statistics,
demonstrating that multinomial classification frameworks can be successfully
applied to Philippine sports data.

\textcite{reyes2021} developed a predictive model for volleyball match
outcomes at the UAAP using ordinal logistic regression, highlighting the
importance of interaction terms in capturing combined effects.

\textcite{mendoza2023} examined the use of multivariate statistical methods
for classifying students into performance tiers at Ateneo de Manila
University, finding that Box's M test should guide the choice between
discriminant analysis and logistic regression.

\textcite{santos2020} studied feature engineering for classification problems
in the Philippine context, emphasizing domain-specific interaction terms over
raw variables.

\subsection{Conceptual Framework}

The study is grounded in the premise that a driver's finishing tier can be
modeled as a function of three categories of predictors:

\begin{enumerate}[label=(\arabic*)]
  \item \textbf{Baseline Pace Control} --- Grid Position captures the
    intrinsic car--driver combination's qualifying performance, absorbing
    the aerodynamic and engine capability differential between teams.
  \item \textbf{Tire Strategy} --- Tire Age (degradation) and Compound
    (Soft, Medium, Hard) represent the strategic choice of rubber and its
    wear trajectory over a stint.
  \item \textbf{Physical Interaction} --- Fuel Mass Proxy (linear burn-off
    from 110\,kg to 0\,kg) and the interaction term Degradation $\times$
    Fuel capture the combined effect of decreasing car mass and increasing
    tire wear---two forces that compound over a race distance.
\end{enumerate}

\noindent
The interaction term is mean-centered prior to multiplication to prevent
structural multicollinearity, allowing the coefficient to be interpreted as
the \textit{combined} effect of both forces rather than their individual
contributions.

\subsection{Research Gap}

The reviewed literature reveals three gaps this study addresses. First,
existing F1 prediction models operate at the season or race-aggregate level
and do not exploit within-race telemetry dynamics. Second, tire degradation
and fuel load are individually recognized as critical race variables, but
their interaction has not been explicitly modeled as a predictor of finishing
tier. Third, most studies validate with random train-test splits rather than
group-aware cross-validation, producing optimistic performance estimates that
do not reflect true generalizability to unseen circuits. This study
synthesizes these gaps by engineering a physically meaningful interaction
term and validating it through LOGO-CV across architecturally diverse tracks.

\subsection{Definition of Terms}

\begin{description}[style=nextline]
  \item[Box's M Test] A multivariate statistical test that evaluates whether
    the variance--covariance matrices of two or more groups are equal.
  \item[Cross-Validation] A model validation technique that partitions data
    into training and testing subsets to assess generalizability.
  \item[Degradation $\times$ Fuel Interaction] The mean-centered product of
    tire age and fuel mass proxy, capturing the combined physical effect of
    rubber wear and decreasing car mass.
  \item[DNF (Did Not Finish)] A classification for drivers who retire from a
    race due to mechanical failure, accident, or other reasons before
    completing the scheduled distance.
  \item[FastF1] An open-source Python library for accessing Formula One
    timing, telemetry, and session data.
  \item[Finishing Tier] A three-class ordinal variable derived from race
    classification: Podium (P1--P3), Point Scorer (P4--P10), No Points
    (P11+).
  \item[Fuel Mass Proxy] An engineered variable modeling fuel burn-off as a
    linear decrease from 110\,kg at race start to 0\,kg at race end.
  \item[Grid Position] A driver's starting position determined by qualifying
    results; used as a control variable for baseline car--driver pace.
  \item[Leave-One-Group-Out Cross-Validation (LOGO-CV)] A validation strategy
    where the model is iteratively trained on $N-1$ groups and tested on the
    held-out group.
  \item[Log-Odds] The natural logarithm of the odds ratio, representing the
    linear predictor in logistic regression.
  \item[Multinomial Logistic Regression] A generalization of binary logistic
    regression for response variables with three or more unordered categories.
  \item[Maximum Likelihood Estimation (MLE)] A method of estimating model
    parameters by finding values that maximize the probability of observing
    the given data.
  \item[Tire Age] The number of laps a tire has been used, measured by the
    \texttt{TyreLife} field in FastF1 telemetry.
  \item[Tire Compound] The type of tire rubber (Soft, Medium, Hard) selected
    for a stint, each offering different performance and durability trade-offs.
  \item[Variance Inflation Factor (VIF)] A measure of how much the variance of
    a regression coefficient is inflated due to multicollinearity with other
    predictors.
\end{description}
